\documentclass[12pt,a4paper]{article}

% -------------------- PACKAGES --------------------
\usepackage[a4paper,margin=1in]{geometry}
\usepackage{amsmath, amssymb, amsthm}
\usepackage{graphicx}
\usepackage{hyperref}
\usepackage{setspace}
\usepackage{enumitem}

\hypersetup{
    colorlinks=true,
    linkcolor=blue,
    urlcolor=blue
}

\onehalfspacing

% -------------------- TITLE --------------------
\title{\textbf{CSE400 -- Fundamentals of Probability in Computing} \\[4pt]
\large Lecture 8: Gaussian Distribution \\[6pt]
\small Prepared from Lecture Slides: L8\_S1\_A.pdf}
\author{}
\date{}

% -------------------- DOCUMENT --------------------
\begin{document}

\maketitle
\hrule
\vspace{1em}

% ==================================================
\section{Foundations: Moments of a Continuous Random Variable}

\subsection{n-th Order Central Moment}

For a continuous random variable $X$ with mean $\mu_X$:

\[
\mathbb{E}[(X-\mu_X)^n]
=
\int (x-\mu_X)^n f_X(x)\,dx
\]

For the discrete case:

\[
\mathbb{E}[(X-\mu_X)^n]
=
\sum_k (x_k-\mu_X)^n p_X(x_k)
\]

% --------------------------------------------------
\subsection{Evaluation for Specific Orders}

\subsubsection*{(i) Zeroth Order ($n=0$)}

\[
\mathbb{E}[(X-\mu_X)^0]
=
\mathbb{E}[1]
=
1
\]

\subsubsection*{(ii) First Order ($n=1$)}

\[
\mathbb{E}[X-\mu_X]
=
\mathbb{E}[X] - \mu_X
\]

Since $\mathbb{E}[X]=\mu_X$,

\[
\mathbb{E}[X-\mu_X]=0
\]

\subsubsection*{(iii) Second Order ($n=2$): Variance}

\[
\sigma_X^2
=
\mathbb{E}[(X-\mu_X)^2]
\]

Expand:

\[
(X-\mu_X)^2
=
X^2 - 2\mu_X X + \mu_X^2
\]

Taking expectation:

\[
\mathbb{E}[X^2]
-
2\mu_X \mathbb{E}[X]
+
\mu_X^2
\]

Since $\mathbb{E}[X]=\mu_X$,

\[
\sigma_X^2
=
\mathbb{E}[X^2] - \mu_X^2
\]

% --------------------------------------------------
\subsection{Higher Order Central Moments}

\subsubsection*{Skewness ($n=3$)}

\[
C_s
=
\frac{\mathbb{E}[(X-\mu_X)^3]}{\sigma_X^3}
\]

\begin{itemize}
\item $C_s > 0$ $\rightarrow$ Right-skewed PDF
\item $C_s < 0$ $\rightarrow$ Left-skewed PDF
\end{itemize}

\subsubsection*{Kurtosis ($n=4$)}

\[
C_k
=
\frac{\mathbb{E}[(X-\mu_X)^4]}{\sigma_X^4}
\]

Large kurtosis indicates a sharp peak near the mean.

% --------------------------------------------------
\subsection{Linearity of Expectation}

For constants $a,b$:

\[
\mathbb{E}[aX+b]
=
a\mathbb{E}[X] + b
\]

For a sum of functions:

\[
\mathbb{E}\left[\sum_{k=1}^{N} g_k(X)\right]
=
\sum_{k=1}^{N} \mathbb{E}[g_k(X)]
\]

Also,

\[
\mathbb{E}[X]
=
\int x f_X(x)\,dx
\]

% ==================================================
\section{Gaussian Random Variable}

\subsection{Definition}

A Gaussian random variable has PDF:

\[
f_X(x)
=
\frac{1}{\sqrt{2\pi\sigma^2}}
\exp\left(
-\frac{(x-m)^2}{2\sigma^2}
\right)
\]

Notation:

\[
X \sim \mathcal{N}(m,\sigma^2)
\]

where:
\begin{itemize}
\item $m$ = mean
\item $\sigma^2$ = variance
\end{itemize}

% ==================================================
\section{Standard Forms}

\subsection{Error Function}

\[
\text{erf}(x)
=
\frac{2}{\sqrt{\pi}}
\int_0^x e^{-t^2} dt
\]

\subsection{Complementary Error Function}

\[
\text{erfc}(x)
=
1 - \text{erf}(x)
=
\frac{2}{\sqrt{\pi}}
\int_x^\infty e^{-t^2} dt
\]

\subsection{$\Phi$-Function (CDF of Standard Normal)}

\[
\Phi(x)
=
\frac{1}{\sqrt{2\pi}}
\int_{-\infty}^{x}
e^{-t^2/2} dt
\]

\subsection{$Q$-Function (Gaussian Tail Function)}

\[
Q(x)
=
\frac{1}{\sqrt{2\pi}}
\int_x^{\infty}
e^{-t^2/2} dt
\]

\[
Q(x)=1-\Phi(x)
\]

% ==================================================
\section{Relation Between $\Phi$ and $Q$}

\subsection{Evaluating the CDF}

\[
F_X(x)
=
\int_{-\infty}^{x}
\frac{1}{\sqrt{2\pi\sigma^2}}
\exp\left(
-\frac{(y-m)^2}{2\sigma^2}
\right) dy
\]

Let

\[
t=\frac{y-m}{\sigma}
\]

Then

\[
F_X(x)
=
\Phi\left(\frac{x-m}{\sigma}\right)
\]

\subsection{Tail Probability}

\[
P(X>x)
=
Q\left(\frac{x-m}{\sigma}\right)
\]

Gay identities:

\[
Q(x)=1-\Phi(x)
\]

\[
F_X(x)
=
1 - Q\left(\frac{x-m}{\sigma}\right)
\]

% ==================================================
\section{Worked Example}

Given:

\[
f_X(x)
=
\frac{1}{\sqrt{8\pi}}
\exp\left(
-\frac{(x+3)^2}{8}
\right)
\]

Comparing with Gaussian form:

\[
m=-3,
\quad
2\sigma^2=8
\Rightarrow
\sigma^2=4
\Rightarrow
\sigma=2
\]

Thus:

\[
X \sim \mathcal{N}(-3,4)
\]

\subsection*{(i) $P(X \le 0)$}

\[
P(X \le 0)
=
\Phi\left(\frac{0-(-3)}{2}\right)
=
\Phi\left(\frac{3}{2}\right)
\]

\subsection*{(ii) $P(X>4)$}

\[
P(X>4)
=
Q\left(\frac{4-(-3)}{2}\right)
=
Q\left(\frac{7}{2}\right)
\]

\subsection*{(iii) $P(|X+3|<2)$}

\[
-2 < X+3 < 2
\]

\[
-5 < X < -1
\]

\[
P(-5<X<-1)
=
\Phi\left(\frac{-1+3}{2}\right)
-
\Phi\left(\frac{-5+3}{2}\right)
\]

\subsection*{(iv) $P(|X-2|>1)$}

\[
X<1 \quad \text{or} \quad X>3
\]

\[
P(X<1) + P(X>3)
\]

% ==================================================
\section{Exam-Oriented Summary}

\[
f_X(x)
=
\frac{1}{\sqrt{2\pi\sigma^2}}
\exp\left(-\frac{(x-m)^2}{2\sigma^2}\right)
\]

\[
Z=\frac{X-m}{\sigma}
\]

\[
F_X(x)=\Phi\left(\frac{x-m}{\sigma}\right)
\]

\[
P(X>x)=Q\left(\frac{x-m}{\sigma}\right)
\]

\[
Q(x)=1-\Phi(x)
\]

\end{document}
